\subsection{Depth along a closed subset}

Let A be a Noetherian ring, F a closed subset of \(\operatorname{Spec}(A)\) and M an A-module. By cor. 2 of prop. 2 of no. 1, the element \(\operatorname{prof}_{A}(J;M)\) of \(\mathbf{N} \cup \{+\infty\}\) does not depend on the ideal J of A such that \(F = V(J)\) ; it is called the depth of M along F and is denoted by \(\operatorname{prof}_{F}(M)\) .

Remarks.-- 1) Let \(r\) be an integer. By prop. 2 of no. 1 and II, § 4, no. 4, cor. 2 of prop. 17, the inequality \(\mathrm{prof}_{\mathrm{F}}(\mathbf{M}) \geqslant r\) is equivalent to the following property: for every finitely generated A-module N whose support is contained in F, we have \(\mathrm{Ext}_{\mathrm{A}}^{i}(\mathrm{N},\mathrm{M}) = 0\) for \(i < r\) .

\begin{enumerate}
\def\labelenumi{\arabic{enumi})}
\setcounter{enumi}{1}
\tightlist
\item
  Suppose that the A-module M is finitely generated. By Remark 2 of No. 1 and Th. 2 of No. 4, we have the following equivalences
\end{enumerate}

\[
\begin{array}{c} \operatorname{prof} _ {\mathrm{F}} (\mathrm{M}) = 0 \Longleftrightarrow \operatorname{Ass} (\mathrm{M}) \cap \mathrm{F} \neq \emptyset \\ \operatorname{prof} _ {\mathrm{F}} (\mathrm{M}) <   + \infty \Longleftrightarrow \operatorname{Supp} (\mathrm{M}) \cap \mathrm{F} \neq \emptyset . \end{array}
\]

PROPOSITION 8. Let A be a Noetherian ring, and let F be a closed subset of \(\operatorname{Spec}(\Lambda)\) and let M be a finitely generated A-module. We have

\[
\operatorname{prof} _ {F} (M) = \inf _ {\mathfrak {p} \in F} \operatorname{prof} _ {A _ {\mathfrak {p}}} \left(M _ {\mathfrak {p}}\right) = \inf _ {\mathfrak {p} \in \operatorname{Supp} (M) \cap F} \operatorname{prof} _ {A _ {\mathfrak {p}}} \left(M _ {\mathfrak {p}}\right).
\]

This is clear if \(\operatorname{prof}_{F}(M) = +\infty\) (remark 2). If \(\operatorname{prof}_{F}(M) = 0\) , there exists a prime ideal \(\mathfrak{p} \in \operatorname{Ass}(M) \cap F\) (remark 2); we have \(\mathfrak{p}A_{\mathfrak{p}} \in \operatorname{Ass}(M_{\mathfrak{p}})\) (IV, § 1, no. 2, prop. 5), hence \(\operatorname{prof}_{\Lambda_{\mathfrak{p}}}(\mathrm{M}_{\mathfrak{p}}) = 0\) (remark 2 of No. 1), whence the proposition in this case.

Suppose \(0 < \operatorname{prof}_{\mathrm{F}}(\mathrm{M}) < +\infty\). Let J be an ideal of A such that \(\mathrm{V}(\mathrm{J}) = \mathrm{F}\) , and let \(x\) an element of J such that the homothety \(x_{\mathrm{M}}\) is injective (loc. cit.). For each prime ideal \(\mathfrak{p}\) the homothety \(x_{\mathrm{M}_{\mathfrak{p}}}\) is injective. By prop. 7 of no. 4, we therefore have

\[
\begin{array}{c} \operatorname{prof} _ {\mathrm{F}} (\mathrm{M} / x \mathrm{M}) = \operatorname{prof} _ {\mathrm{F}} (\mathrm{M}) - 1 \\ \operatorname{prof} _ {\Lambda_ {\mathfrak {p}}} \big ((\mathrm{M} / x \mathrm{M}) _ {\mathfrak {p}} \big) = \operatorname{prof} _ {\Lambda_ {\mathfrak {p}}} (\mathrm{M} _ {\mathfrak {p}}) - 1. \end{array}
\]

We then conclude by induction on the integer \(\mathrm{prof}_{\mathbb{F}}(\mathbb{M})\) .

Remark 3. If q is a point of Supp(M), we therefore have \(\operatorname{prof}_{\mathrm{A}}(\mathfrak{q};\mathrm{M}) = \inf_{\mathfrak{p}\supseteq q} \operatorname{prof}_{\Lambda_{\mathfrak{p}}}(\mathrm{M}_{\mathfrak{p}})\) In particular, we have the inequality \(\operatorname{prof}_{\mathrm{A}}(\mathfrak{q};\mathrm{M}) \leqslant \operatorname{prof}_{\mathrm{A}_{\mathfrak{q}}}(\mathrm{M}_{\mathfrak{q}})\) ; equality holds when q is maximal. In the general case, one may have \(\operatorname{prof}_{\mathrm{A}}(\mathfrak{q};\mathrm{M}) < \operatorname{prof}_{\Lambda_{\mathfrak{q}}}(\mathrm{M}_{\mathfrak{q}})\) ; one may also have \(\operatorname{prof}_{\mathrm{A}}(\mathfrak{q};\mathrm{M}) < \inf \operatorname{prof}_{\mathrm{A}_{\mathfrak{m}}}(\mathrm{M}_{\mathfrak{m}})\) where m ranges over the set of maximal ideals of A containing q. Let, for example, p be a nonmaximal prime ideal of A, containing q and distinct from q; set M = A/p. One has \(\operatorname{prof}_{\mathrm{A}}(\mathfrak{q};\mathrm{M}) = 0\) , \(\operatorname{prof}_{\mathrm{A}_{\mathfrak{q}}}(\mathrm{M}_{\mathfrak{q}}) = +\infty\) and \(\operatorname{prof}_{\mathrm{A}_{\mathfrak{m}}}(\mathrm{M}_{\mathfrak{m}}) > 0\) for every maximal ideal m of A.

PROPOSITION 9.--- Let A be a Noetherian ring, M and N two finitely generated A-modules, and F the support of N. Then \(\operatorname{prof}_{\mathbf{F}}(\mathbf{M})\) is the lower bound (in \(\mathbf{N} \cup \{+\infty\}\) ) of the set of integers \(n\) such that \(\operatorname{Ext}_{\Lambda}^{n}(\mathbf{N}, \mathbf{M}) \neq 0\) .

By Remark 1, we have \(\mathrm{Ext}_{\mathrm{A}}^{i}(\mathrm{N},\mathrm{M}) = 0\) for every \(i < \mathrm{prof}_{\mathrm{F}}(\mathrm{M})\) . It remains to prove that if \(\mathrm{prof}_{\mathrm{F}}(\mathrm{M}) = n < +\infty\) , one has \(\mathrm{Ext}_{\mathrm{A}}^{n}(\mathrm{N},\mathrm{M}) \neq 0\) . Let J be the annihilator of N; we have \(\mathrm{F} = \mathrm{V}(\mathrm{J})\) , hence \(\mathrm{prof}_{\mathrm{F}}(\mathrm{M}) = \mathrm{prof}_{\mathrm{A}}(\mathrm{J};\mathrm{M})\) . By cor. 1 of th. 2 ( \(n^{\circ}4\) ), there exists an M-regular sequence ( \(x_{1},\ldots,x_{n}\) ) of length \(n\) consisting of elements of J, and the depth of the A-module relative to J \(\overline{\mathrm{M}} = \mathrm{M}/(x_{1}\mathrm{M}+\ldots+x_{n}\mathrm{M})\) is zero. By A, X, p. 166, prop. 9, it suffices to prove that \(\mathrm{Hom}_{\mathrm{A}}(\mathrm{N},\overline{\mathrm{M}})\) is nonzero. Now, by prop. 8, there exists \(\mathfrak{p} \in \mathrm{Supp}(\mathrm{M}) \cap \mathrm{Supp}(\mathrm{N})\) such that \(\mathrm{prof}_{\mathrm{A}_{\mathfrak{p}}}(\overline{\mathrm{M}}_{\mathfrak{p}}) = 0\) , that is \(\mathrm{Hom}_{\mathrm{A}_{\mathfrak{p}}}(\kappa(\mathfrak{p}),\overline{\mathrm{M}}_{\mathfrak{p}}) \neq 0\) . Since \(\mathbf{N}_{\mathfrak{p}}\) is nonzero, the \(\kappa(\mathfrak{p})\) -vector space \(\mathbf{N}_{\mathfrak{p}} \otimes_{\mathbf{A}_{\mathfrak{p}}} \kappa(\mathfrak{p})\) is nonzero (Nakayama's lemma), as is its dual; there therefore exists a surjective \(\mathbf{A}_{\mathfrak{p}}\) -linear map from \(\mathbf{N}_{\mathfrak{p}}\) in \(\kappa(\mathfrak{p})\) . It follows that one has \(\mathrm{Hom}_{\mathrm{A}_{\mathfrak{p}}}(\mathbf{N}_{\mathfrak{p}},\overline{\mathbf{M}}_{\mathfrak{p}}) \neq 0\) , hence \(\mathrm{Hom}_{\mathrm{A}}(\mathrm{N},\overline{\mathrm{M}}) \neq 0\) (II, § 2, \(n^{\circ}7\) , prop. 19), which is what we wanted to prove.

Remark 4.--- Let A be a Noetherian ring, N a finitely generated A-module. One sometimes calls the grade of N, and denotes it by grade (N), the greatest lower bound in \(\mathbf{N}\cup\{+\infty\}\) of the set of integers \(n\) such that \(\operatorname{Ext}_{\mathrm{A}}^{n}(\mathrm{N},\mathrm{A})\) is nonzero. By prop. 9, this is also the depth of A along the support of N, or again ( \(n^{\circ}\) 4, th. 2) the least upper bound of the set of lengths of A-regular sequences of elements of the annihilator of N. Since for every prime ideal p of A the annihilator of \(N_{p}\) is equal to \(\operatorname{Ann}(N)_{p}\) (II, § 2, \(n^{\circ}\) 4, formula (9)), one deduces from prop. 5 of \(n^{\circ}\) 3 the equality

\[
\operatorname{grade} (\mathrm{N}) = \inf _ {\mathfrak {p} \in \operatorname{Spec} (\Lambda)} \operatorname{grade} \left(\mathrm{N} _ {\mathfrak {p}}\right) = \inf _ {\mathfrak {m} \in \Omega} \operatorname{grade} \left(\mathrm{N} _ {\mathfrak {m}}\right),
\]

where Ω denotes the set of maximal ideals of A.

